\documentclass[11pt,a4paper]{article}
\usepackage[margin=1in]{geometry}
\usepackage{enumitem}
\usepackage{booktabs}
\usepackage{tabularx}
\usepackage{hyperref}
\usepackage{amsmath}
\usepackage{titling}
\usepackage{graphicx}
\usepackage{longtable}
\usepackage{xcolor}
\usepackage{ragged2e}
\emergencystretch=2em

% -----------------------------------------------------
% bvraghav-tiet-ucs503p-article.sty (as used in the proposal)
% -----------------------------------------------------
\makeatletter \newcommand{\BvrSetLabInstructor}[1]{%
  \gdef\Bvr@LabInstructor{#1}}
\newcommand{\Bvr@LabInstructor}{Dr. Raghav %
  B. Venkataramaiyer}
\newcommand{\BvrLabInstructorValue}{\Bvr@LabInstructor}
\makeatother
\pretitle{\begin{center}\bfseries\fontsize{18bp}{18bp}\selectfont}
\makeatletter
\newcommand{\BvrSubtitle}[1]{%
  \posttitle{%
    \par\end{center}
    \begin{center}\large#1\end{center}
    \vskip 0.5em}%
}
\makeatother

\newcommand{\BvrHint}[1]{%
  {\normalfont\normalsize\itshape\hfill #1}}

\newcommand{\rollno}[1]
{Roll No: \texttt{#1}}
\newcommand{\email}[1]
{Email: \texttt{\href{mailto:#1}{#1}}}
\newcommand{\specifyauthor}[3]
{\texttt{#1}, CSED\thanks{\rollno{#2}, \email{#3}}}

% -----------------------------------------------------

\title{SCBE (Semi Computer Based Evaluation)}
\BvrSetLabInstructor{Dr. Raghav B. Venkataramaiyer}

\BvrSubtitle{%
  Project Report --- Prototype Stage, for UCS503P \\
  Submitted to: \BvrLabInstructorValue \\ \bfseries
  Thapar Institute of Engineering and Technology%
}

\renewcommand{\maketitlehookb}{%
  \begin{center}
    \large\textbf{Team Infinix}\par
    \vspace{0.5em}
  \end{center}
}

\setlength{\skip\footins}{4em}

\author{
  \specifyauthor{Gurneet Singh Ahuja}{1024170014}{gahuja\_be24@thapar.edu}
  \and
  \specifyauthor{Alla Harshith}{1024170007}{aharshith\_be24@thapar.edu}
  \and
  \specifyauthor{Arpit Goyal}{1024170021}{agoyal6\_be24@thapar.edu}
}

\date{\today}
\begin{document}
\maketitle
\tableofcontents
\vspace{1em}

\section{Introduction}

SCBE (Semi Computer Based Evaluation) is a platform for digitizing the post-examination
answer-sheet review process. In the workflow the platform targets, faculty capture a
student's physical answer booklet as a short video or a set of photographs; the platform
assembles these captures into a single PDF, lets faculty attach the sheet to the correct
student and examination, and publishes it so the student can view and annotate their own
evaluated answer sheet remotely. If a student disagrees with part of the evaluation, they
can raise a complaint directly on the annotated page, and the faculty member can review and
resolve it from their own dashboard. The proposal framed this as a way to cut down on the
travel, delay, and lack of a paper trail that come with the existing practice of physically
re-checking answer sheets on campus.

Since the proposal, the team has built a working three-tier prototype -- a Python/OpenCV
service that turns raw video or images into a paginated PDF, a Node.js/Express backend that
owns authentication, data persistence and orchestration, and a React frontend with
role-specific dashboards for students, faculty, and an examination-cell admin. This report
describes what has actually been implemented, how the pieces fit together, and what has
changed relative to the original plan.

\subsection{Objectives: status at the prototype stage}

\begin{longtable}{>{\RaggedRight}p{0.40\textwidth} >{\RaggedRight}p{0.12\textwidth} >{\RaggedRight}p{0.40\textwidth}}
\toprule
\textbf{Objective (from proposal)} & \textbf{Status} & \textbf{Notes} \\
\midrule\endhead
Faculty and student authentication with access controls & Met &
  Role-based login/registration (student, faculty, admin) with salted, hashed
  passwords; routes are scoped per role and per record owner. \\[0.6em]
Faculty workflow for uploading and managing evaluated answer sheets & Met &
  Video-to-PDF and images-to-PDF conversion, plus page-level add/delete/replace
  editing, are implemented end to end. \\[0.6em]
Student interface for accessing and viewing evaluated sheets & Met &
  In-browser PDF viewer with per-question mark totals and evaluation
  breakdown. \\[0.6em]
Student workflow for raising review requests & Met &
  Students annotate directly on the PDF and submit a complaint tied to a
  specific question and the annotated page. \\[0.6em]
Faculty workflow for reviewing and responding to requests & Met &
  Faculty can inspect the annotated PDF, change complaint status
  (Pending / Under Review / Resolved / Rejected), and attach a comment. \\[0.6em]
Recording of requests, responses, and resolutions & Met &
  Persisted in a relational schema (complaints, discrepancies, notifications)
  with an audit trail of status changes. \\[0.6em]
CI/CD with automated testing to a hosted staging environment & Not started &
  Deferred; the prototype currently runs from local processes only. See
  Section~\ref{sec:limitations}. \\[0.6em]
Automatic association of answer sheets with the correct student (OCR) & Partially met &
  The database and API already carry \texttt{detectedRollNumber} /
  \texttt{detectedName} fields, but the OCR/handwriting-recognition step that
  should populate them has not been built yet; faculty currently assign each
  sheet to a student manually. \\[0.6em]
Faculty verification/correction of automatic associations & Deferred &
  Depends on the OCR step above; the manual assignment screen that will later
  serve as the ``correction'' step already exists. \\
\bottomrule
\end{longtable}

The most significant deviation from the proposal is on automatic
answer-sheet identification. The proposal treated OCR-based name/roll-number
extraction as part of the initial deliverables; in practice the team
prioritized getting the full manual pipeline -- capture, convert, assign,
publish, review, complain, resolve -- working reliably first, and pushed OCR
to the next iteration. This is discussed further in Section~\ref{sec:challenges}.

\section{System Architecture}

The prototype is organised as three independently runnable services that
communicate over HTTP, matching the decoupled, stateless design the proposal
argued for on scalability grounds.

\begin{center}
\texttt{React frontend (Vite)} $\;\longrightarrow\;$
\texttt{Node.js/Express backend (port 4000)} $\;\longrightarrow\;$
\texttt{Flask/OpenCV CV service (port 5000)}
\end{center}

\begin{itemize}[leftmargin=0.7cm]
  \item \textbf{Frontend (React + Vite):} role-specific pages for students,
    faculty, and admins, an in-browser PDF viewer/annotator
    (\texttt{@embedpdf/react-pdf-viewer}), and a thin \texttt{api.js} client
    that talks to the backend.
  \item \textbf{Backend (Node.js + Express):} authentication, the
    \texttt{users}/\texttt{courses}/\texttt{exams}/\texttt{sheets}/\texttt{complaints}/
    \texttt{discrepancies}/\texttt{notifications} data model, file storage for
    uploaded and generated documents, and a thin proxy layer that forwards
    capture jobs to the CV service.
  \item \textbf{CV service (Python + Flask + OpenCV):} stateless
    image/video-processing routes that turn a video or a set of images into a
    PDF, and that perform page-level edits (delete/add/replace) on an existing
    PDF.
  \item \textbf{Storage:} a single SQLite database (\texttt{scbe.sqlite}, in
    WAL mode) for structured data, and a directory tree on disk for uploaded
    videos/images and generated original/annotated PDFs.
\end{itemize}

\subsection{Core workflow as implemented}

\begin{enumerate}[leftmargin=0.7cm]
  \item Faculty record a short video (or take photos) of the evaluated answer
    booklet and upload it from the \emph{Upload} screen.
  \item The backend streams the file to the CV service, which returns a
    generated PDF; the backend stores the PDF and creates an unassigned
    \texttt{sheet} row.
  \item From the \emph{Assign} screen, faculty attach the sheet to a
    student and an exam (roll number/name can currently be entered manually;
    see Section~\ref{sec:limitations} on OCR).
  \item Faculty publish the sheet, which triggers a notification to the
    student and makes the sheet visible on their dashboard.
  \item The student opens the sheet in the PDF viewer, sees the per-question
    marks and total, and can annotate a page and raise a complaint about a
    specific question.
  \item The complaint (with the flattened annotated PDF and the raw
    annotation JSON) lands on the faculty's complaint queue; the faculty
    member reviews it in the same viewer, changes its status, and adds a
    resolution comment, which the student can then see.
\end{enumerate}

\subsection{Technical stack}

\begin{table}[h!]
\centering
\begin{tabular}{p{0.28\textwidth} p{0.62\textwidth}}
\toprule
\textbf{Component} & \textbf{Technology} \\
\midrule
Frontend & React 18, React Router, Vite \\
PDF viewing/annotation & \texttt{@embedpdf/react-pdf-viewer} (PDFium/WASM) \\
Backend & Node.js, Express, Multer (file uploads) \\
Database & SQLite (\texttt{node:sqlite}, WAL mode) \\
Password storage & \texttt{scrypt} key derivation with per-user salt \\
CV / document assembly & Python, Flask, Flask-CORS \\
Frame processing & OpenCV (\texttt{opencv-python-headless}), NumPy \\
Image/PDF assembly & Pillow, PyPDF2 \\
\bottomrule
\end{tabular}
\caption{Technical stack used in the prototype}
\end{table}

\section{Module-by-Module Implementation}

\subsection{Video/image-to-PDF conversion}

This is the module that most directly replaces the proposal's ``dedicated
assistant scans the sheets'' step. Rather than requiring faculty to pick out
individual page photographs by hand, the CV service accepts a single video of
the whole booklet being flipped through and extracts one clean frame per
page.

The extraction logic samples the incoming video at a fixed rate
(default 20 fps) and tracks a simple two-state machine per sampled frame:

\begin{itemize}[leftmargin=0.7cm]
  \item \textbf{STABLE:} the camera/page is holding still. Every sampled frame
    is kept as a candidate for the current page, scored by a Laplacian-variance
    sharpness measure (a common focus-quality heuristic: a sharper image has
    more high-frequency edge content, hence higher Laplacian variance).
  \item \textbf{FLIPPING:} consecutive frame-to-frame differences (mean
    absolute pixel difference on grayscale frames) exceed a threshold for
    enough consecutive samples, which is taken as evidence that a page turn is
    in progress.
\end{itemize}

When the frame difference stays low for long enough after a flip, the state
returns to STABLE and a new page begins accumulating candidates. When a flip
is detected, the sharpest candidate frame collected during the previous
STABLE run is committed as a finished page. All committed frames are finally
stitched, in order, into a single multi-page PDF. Sample rate, the
difference threshold, and the number of consecutive frames required to
confirm a flip or a stabilization are all exposed as tunable parameters on
the API rather than hard-coded, which was useful during testing under
different lighting and camera-shake conditions.

The same service also accepts a batch of individual page images directly
(for faculty who prefer photographing pages one at a time) and stitches them
into a PDF in upload order, and it can delete, insert, or replace a single
page of an already-generated PDF -- used when a page was mis-captured and
needs to be redone without regenerating the whole document.

\subsection{Backend orchestration and data model}

The Node backend does not perform any image processing itself; it is
responsible for authentication, persistence, and gluing the frontend to the
CV service. Uploaded video/image files are received with Multer, forwarded
to the CV service as \texttt{multipart/form-data}, and the resulting PDF is
copied into a durable \texttt{generated/pdfs/originals} directory referenced
by a new \texttt{sheets} row.

The schema (SQLite) currently has seven tables: \texttt{users},
\texttt{courses}, \texttt{exams}, \texttt{sheets}, \texttt{discrepancies},
\texttt{notifications}, \texttt{complaints}, and a small \texttt{meta} table
used for lightweight schema-version bookkeeping. The \texttt{sheets} table
already carries \texttt{detectedRollNumber} and \texttt{detectedName}
columns, reflecting the planned OCR integration, even though nothing writes
to them automatically yet -- they are filled in by faculty through the
manual assignment endpoint. Passwords are never stored in plaintext:
each user gets a random salt, and the derived \texttt{scrypt} hash is
compared using a constant-time comparison to reduce timing side-channels.

\subsection{Annotation and complaint capture}

Students view their sheet using \texttt{@embedpdf/react-pdf-viewer}, which
runs PDF rendering through a PDFium WebAssembly build in the browser rather
than round-tripping page images through the server. When a student opens the
complaint form on a page, the viewer's annotation plugin exports the
annotation set for that document (through its \texttt{exportAnnotations}
capability); the frontend attaches this JSON alongside the already-flattened,
annotated PDF and the complaint title/description, and posts it as a single
multipart request. This gives the faculty member two views of the same
concern when they open the complaint later: the flattened PDF (a durable
record that cannot be re-edited after submission) and the structured
annotation data (useful if the interface ever needs to re-render annotations
interactively instead of only showing the flattened image).

\subsection{Faculty and student dashboards}

Faculty have four working screens: a dashboard summarising counts and recent
resolved complaints, a sheets list (assigned/unassigned, with per-sheet
publish and reassignment actions), the capture/upload screen described
above, and a complaint inbox with status transitions
(\emph{Pending} $\to$ \emph{Under Review} $\to$ \emph{Resolved}/\emph{Rejected}).
Students have a matching set: dashboard, sheet list, the annotated PDF
viewer with the complaint form, a list of their own submitted complaints
with status, and notifications. The admin role currently exposes read-only
aggregate views (student/faculty counts, examinations, and a global list of
answer sheets) for the examination cell to get a bird's-eye view of platform
activity; write operations for admins (e.g. course/exam management from the
UI) are not yet implemented.

\section{Testing and Validation}
\label{sec:testing}

Formal automated CI/CD and a load-tested staging deployment were listed in
the proposal as initial deliverables but have not been reached at this
stage; testing so far has been manual and workflow-oriented rather than
based on an automated test suite. The team exercised the following
scenarios directly against the running services:

\begin{itemize}[leftmargin=0.7cm]
  \item \textbf{End-to-end happy path:} register a faculty and a student
    account, upload a sample video of a mock answer booklet, confirm a PDF is
    generated, assign it to the student, publish it, and confirm it appears
    on the student's dashboard.
  \item \textbf{Page-extraction quality:} varied the sample video's flip
    speed and lighting to check that the difference-threshold/sharpness
    heuristic still picks one legible frame per page; very fast flips or low
    light occasionally required raising \texttt{sample\_fps} or lowering
    \texttt{difference\_threshold} for reliable page separation.
  \item \textbf{PDF page editing:} deleted, inserted, and replaced a page on
    an already-generated PDF and confirmed the page count and page order in
    the output.
  \item \textbf{Complaint round-trip:} annotated a sheet as a student,
    submitted a complaint, confirmed the faculty inbox received it with the
    flattened PDF intact, changed its status with a comment, and confirmed
    the status/comment appeared back on the student side.
  \item \textbf{Access control:} confirmed a student cannot open another
    student's sheet or complain about a sheet that is not theirs, and that
    role-scoped routes reject requests for a user id/role that do not match
    the record (e.g. a student id used against a faculty-only route).
\end{itemize}

No formal performance benchmarking (throughput under concurrent uploads,
response time under load) has been carried out yet; the proposal's
``prototype workload testing'' and ``document delivery performance'' items
under Practical Scalability remain future work, as does the hosted
deployment needed to test under realistic network conditions.

\subsection{Result: what currently works vs.\ what does not}

\begin{longtable}{>{\RaggedRight}p{0.44\textwidth} >{\RaggedRight}p{0.12\textwidth} >{\RaggedRight}p{0.36\textwidth}}
\toprule
\textbf{Capability} & \textbf{Result} & \textbf{Notes} \\
\midrule\endhead
Video $\to$ PDF conversion & Pass &
  Reliable for moderate flip speed and adequate lighting. \\[0.4em]
Images $\to$ PDF conversion & Pass &
  Direct, order-preserving stitching. \\[0.4em]
PDF page delete/add/replace & Pass &
  Verified against multi-page sample PDFs. \\[0.4em]
Manual sheet assignment \& publish & Pass &
  Notification is created and appears for the student. \\[0.4em]
Student annotation \& complaint submission & Pass &
  Flattened PDF and annotation JSON both persisted. \\[0.4em]
Faculty complaint review/resolution & Pass &
  Status and resolution comment both round-trip to the student view. \\[0.4em]
Per-student/per-role access control & Pass (spot-checked) &
  No automated authorization test suite yet. \\[0.4em]
Automatic student identification (OCR) & Not implemented &
  Field exists in schema; extraction pipeline not built. \\[0.4em]
CI/CD and hosted deployment & Not implemented &
  Prototype currently runs from local processes only. \\
\bottomrule
\end{longtable}

\section{Challenges and Mitigations}
\label{sec:challenges}

\begin{itemize}[leftmargin=0.7cm]
  \item \textbf{Separating a genuine page turn from camera shake.} An early
    version of the video-processing heuristic mistook ordinary camera wobble
    for a page flip, which fragmented single pages into several spurious
    ``pages''. Requiring several consecutive frames to cross the difference
    threshold before committing to a flip (rather than reacting to a single
    frame) removed most of the false positives without needing a heavier
    motion-tracking model.
  \item \textbf{Choosing a representative frame per page.} Even within a
    stable interval, individual frames vary in focus. Scoring every candidate
    frame with a Laplacian-sharpness measure and keeping the sharpest one per
    page gave a consistent improvement over simply taking the first or last
    frame of the interval.
  \item \textbf{Deferring OCR-based identification.} The proposal's
    ``automatic answer-sheet identification'' component turned out to need
    more groundwork (a labelled dataset of the institute's roll-number/name
    formats, and a decision on which handwriting-recognition engine to
    integrate) than the other, more mechanical parts of the pipeline. Rather
    than block the rest of the prototype on this, the team built the manual
    assignment path first and designed the schema
    (\texttt{detectedRollNumber}, \texttt{detectedName}) so that an OCR step
    can later populate the same fields the manual path already writes to,
    with faculty verification staying in the loop exactly as the proposal's
    fallback-handling plan intended.
  \item \textbf{Keeping the annotated record trustworthy.} Because a
    complaint doubles as an integrity record of what the student is disputing,
    the frontend submits a flattened, already-rendered PDF snapshot rather
    than only the editable annotation JSON, so a later change to the viewer
    or a client-side bug cannot silently alter what was submitted at
    complaint time.
\end{itemize}

\section{Limitations}
\label{sec:limitations}

\begin{itemize}[leftmargin=0.7cm]
  \item Answer-sheet–to–student association is manual; there is no automatic
    OCR/handwriting-recognition step yet, so the ``identification accuracy''
    metric from the proposal cannot be measured at this stage.
  \item There is no CI/CD pipeline or hosted staging deployment; the system
    has only been run and tested locally.
  \item No automated test suite (unit or integration) exists yet; validation
    so far has been manual, scenario-based testing as described in
    Section~\ref{sec:testing}.
  \item The admin role is read-only in the current UI.
  \item Performance/scalability has not been measured under concurrent load
    or with large batches of sheets.
\end{itemize}

\section{Conclusion and Future Work}

The prototype demonstrates a complete, working loop for the core workflow the
proposal set out to build: faculty can turn a recorded video of a physical
answer booklet into a digital PDF, attach it to the right student, publish
it, and handle any resulting dispute through an auditable
annotate-complain-resolve cycle, all without a student needing to visit
campus. The main piece of the original proposal not yet delivered is
automatic student identification from the scanned sheet; everything around
it -- the schema fields, the manual verification/override screen, and the
downstream publish/notify pipeline -- is already in place to receive it.

Planned next steps, in rough priority order, are:

\begin{enumerate}[leftmargin=0.7cm]
  \item Integrate a handwriting-recognition/OCR step (e.g.\ Tesseract,
    PaddleOCR, or TrOCR, as originally scoped) to populate
    \texttt{detectedRollNumber}/\texttt{detectedName} automatically, with the
    existing manual-assignment screen repurposed as the faculty
    verification/correction step.
  \item Add automated tests (unit tests for the CV heuristics and page-edit
    operations, integration tests for the core API routes) and a basic CI
    pipeline.
  \item Move from local processes to a hosted staging deployment, and run the
    pilot validation plan described in the proposal (10--20 students,
    timed comparison against the physical process).
  \item Benchmark the CV service and backend under concurrent uploads and
    concurrent PDF downloads to establish real throughput numbers ahead of
    the proposal's scalability claims.
  \item Extend the admin role beyond read-only dashboards (course/exam
    management, user administration).
\end{enumerate}

\end{document}
